\documentclass[11pt,a4paper]{article}
\usepackage[utf8]{inputenc}
\usepackage[top=1.9cm,bottom=1.9cm,left=2.1cm,right=2.1cm]{geometry}
\usepackage{booktabs}
\usepackage{titlesec}
\usepackage{parskip}

\titleformat{\section}{\normalfont\large\bfseries}{\thesection.}{0.5em}{}
\titlespacing*{\section}{0pt}{10pt}{4pt}

\newcommand{\entry}[1]{\par\hspace*{1.5em}#1}

\begin{document}

\begin{center}
    {\large \textbf{CAPSTONE PROJECT -- TOPIC DESCRIPTION}}\\[4pt]
    {\large \textbf{A Web-Based System for Real-Time Hotel Market Price
    Monitoring and Forecasting Using Machine Learning}}\\[6pt]
    Student: Dang Nguyen Minh An \hfill Degree: Master of Software Engineering
\end{center}

\vspace{-2pt}
\hrule
\vspace{6pt}

\section{Topic Orientation and Motivation}
Hotel room prices on Online Travel Agencies (OTAs) such as Agoda and
Booking.com fluctuate continuously under dynamic pricing driven by demand,
seasonality, events, and competitor behavior. Large hotel chains manage this
volatility with commercial revenue management and rate intelligence platforms
(e.g., RateGain, Lighthouse, Duetto), which cost hundreds to thousands of USD
per month. However, Vietnam's accommodation sector - over 38{,}000
establishments, the majority independent small and medium-sized hotels -
largely cannot access such tools: competitor rates are checked manually and
room prices are set by intuition. On the other side of the market, individual
travelers face the mirror problem of price opacity, with no way to track
price history or anticipate price movements.

This project develops an end-to-end web-based market price intelligence
system for the Vietnamese hotel market that: (1)~continuously collects room
prices from Agoda and Booking.com across major Vietnamese cities,
(2)~applies machine learning to forecast short-term price movements (1--14
days ahead), and (3)~serves the forecasts through two decision layers built
on a shared prediction engine:

\entry{\textbf{Primary use case - Hotelier / Business view:} a competitor
rate intelligence dashboard where a hotel selects its competitive set
(comparable nearby hotels), monitors their real-time rates across both OTAs,
views market-level price forecasts, receives alerts when its own price
deviates significantly from the market or when market movements are
predicted (holidays, events), and detects cross-OTA rate disparities.}

\entry{\textbf{Secondary use case - Consumer view:} price history charts and
short-term forecasts for individual hotels, supporting travelers'
booking-timing decisions.}

Both views are powered by a single trained forecasting model; only the
application layer differs. This addresses a clear research and market gap:
affordable, ML-based rate intelligence for the under-explored Vietnamese
market, accounting for local demand drivers (Tet, summer domestic peak,
regional festivals) that existing Western-oriented tools ignore.

\section{Dataset and Data Collection (Current Progress)}
The dataset is self-collected. A web-based crawling tool has already been
implemented and is operational. The tool manages lists of hotels organized by
region/city; for each hotel and each selected check-in/check-out date pair,
it automatically extracts all publicly listed information. Each record
contains:

\entry{\textbf{Price data:} nightly rate (VND), discount, original price, room
type, cancellation policy, breakfast inclusion.}

\entry{\textbf{Hotel attributes:} name, star rating, location (city/district),
review score, number of reviews, amenities.}

\entry{\textbf{Temporal metadata:} scrape timestamp, check-in date, and lead
time (days between observation and check-in) - the key dimension for
modeling how prices evolve as the stay date approaches.}

The crawler can run automatically at predefined times or be triggered
manually on demand by the user. Each run observes a hotel's prices for
multiple future check-in dates, producing dense price time series. Data is
enriched with external public sources: the Vietnamese holiday/festival
calendar, weather forecasts (Open-Meteo API), and VNAT tourism statistics.
Target scale: 200--300 hotels (3--5 stars) in 5 cities (Ho Chi Minh City,
Hanoi, Da Nang, Nha Trang, Phu Quoc), accumulating 100{,}000+ price records
over the collection period. Supervised-learning labels are derived
automatically from the time series itself (future prices are observed later
by the same crawler), so no manual annotation is required.

\section{Methodology and Models}
\entry{\textbf{Feature engineering:} Four feature groups are constructed: hotel
attributes; temporal features (day of week, holiday/festival flags, lead
time, season); price-history features (lag values, rolling mean/std, price
velocity); and market context (competitive-set average price, cross-OTA
price ratio, weather, tourist arrivals).}

\entry{\textbf{Models:} Four models - two traditional machine learning models and
two deep learning models - are trained and systematically compared for
price forecasting at 1, 3, 7, and 14-day horizons. No model is designated in
advance as the primary candidate; the best-performing model is selected
empirically based on experimental results.}

\begin{center}
\begin{tabular}{lll}
\toprule
Model & Family & Paradigm represented \\
\midrule
XGBoost & Traditional ML & Boosting (sequential error correction) \\
Random Forest & Traditional ML & Bagging (variance-reducing ensemble) \\
LSTM & Deep Learning & Recurrent sequence modeling (gated memory) \\
Transformer & Deep Learning & Attention-based sequence modeling \\
\bottomrule
\end{tabular}
\end{center}

\entry{\textbf{Model selection rationale:} Each model represents a methodologically
distinct paradigm, making the comparison meaningful rather than redundant.
XGBoost represents gradient boosting, the family that consistently dominates
tabular prediction benchmarks and prior price-forecasting studies; Random
Forest represents bagging, an independent ensemble principle (variance
reduction rather than bias reduction). LSTM is the canonical recurrent
architecture for temporal dependencies in price series, well-established in
tourism demand forecasting literature, while the Transformer represents the
current attention-based direction in time-series forecasting - providing a
recurrent-vs-attention contrast. Near-identical variants within a family
(LightGBM/CatBoost vs. XGBoost; GRU vs. LSTM) are deliberately excluded:
their differences are mainly engineering optimizations, adding training cost
without methodological insight. Classical statistical models (ARIMA,
exponential smoothing) are excluded because they cannot naturally exploit
the rich exogenous features central to this problem (hotel attributes,
holidays, weather, competitive-set context).}

A strict time-based train/validation/test split (70/15/15, chronological)
prevents data leakage; hyperparameters are tuned with Optuna; SHAP analysis
provides interpretability (which factors drive market price movements -
directly valuable insight for hoteliers). Models are evaluated on both price
regression and price-direction classification (increase/decrease). The
winning model becomes the shared forecasting engine: hotel-level predictions
are served directly to the consumer view and aggregated over competitive
sets for the hotelier view.

\entry{\textbf{System:} Scrapy/Selenium crawlers with scheduled and on-demand
execution, PostgreSQL + Redis storage, FastAPI backend, and a ReactJS
dashboard (competitor rate monitoring, market forecast panel, pricing
position alerts, price history charts).}

\section{Evaluation and Expected Results}
Evaluation is fully quantitative, on held-out chronological test data:
(1)~Forecast accuracy - target MAPE $\leq 12\%$ for 7-day forecasts,
price-direction accuracy $\geq 80\%$, with SHAP-based interpretability
analysis and per-city/per-star-rating breakdowns. (2)~Practical utility via
retrospective simulation - a pricing opportunity analysis for the hotelier
view (quantifying cases where a hotel's listed price deviated from the
predicted market movement during high-demand periods, i.e., potential
mispricing detected by the system), and a booking savings simulation for the
consumer view (savings achieved by following forecast-based timing versus
immediate booking). Expected deliverables: the curated Vietnamese hotel
price dataset, four trained and compared ML models with full experimental
analysis, and a deployable dual-view web application.

\end{document}